\clearpage
\VerdierPhysicalPageAnchor{0246}%
\begingroup\small\setlength{\parskip}{0pt}\setlength{\emergencystretch}{1.5em}
\setlength{\abovedisplayskip}{4pt plus 2pt minus 2pt}%
\setlength{\belowdisplayskip}{4pt plus 2pt minus 2pt}%

\VerdierSemanticAnchor{ch4d-4-6-4}
\noindent\textbf{4.6.4.}\quad
$X^{\bullet\bullet}$ を $\mathcal A$ の二重複体とする（第 I 章、2.1.2）。
二重複体 $X^{\bullet\bullet}$ は
$\operatorname{comp}(\operatorname{comp}(\mathcal A))$ の対象を定める：

\VerdierSemanticAnchor{ch4d-4-6-4-double-complex}
\[
X^{\bullet\bullet}\colon\quad
\cdots\longrightarrow X^{\bullet,n-1}
\xrightarrow{d_2^{\bullet,n-1}}X^{\bullet,n}
\xrightarrow{d_2^{\bullet,n}}X^{\bullet,n+1}
\longrightarrow\cdots.
\]
（(4.1.9) の記法を用いる。符号に注意せよ！）

\VerdierSemanticAnchor{ch4d-4-6-4-b1}
$B_1^p(X^{\bullet\bullet})$ を次の単純複体と記す：

\VerdierSemanticAnchor{ch4d-4-6-4-1}
\[
\left\{
\begin{aligned}
B_1^p(X^{\bullet\bullet})^i
  &=\operatorname{Im}(d_1^{p-1,i}),\\
d_{B_1^p(X^{\bullet\bullet})}^{i}
  &\text{次の微分から誘導される： }d_2^{p,i}.
\end{aligned}
\right.
\tag{4.6.4.1}
\]

\VerdierSemanticAnchor{ch4d-4-6-4-z1}
$Z_1^p(X^{\bullet\bullet})$ を次の単純複体と記す：

\VerdierSemanticAnchor{ch4d-4-6-4-2}
\[
\left\{
\begin{aligned}
Z_1^p(X^{\bullet\bullet})^i
  &=\operatorname{Ker}(d_1^{p,i}),\\
d_{Z_1^p(X^{\bullet\bullet})}^{i}
  &\text{次の微分から誘導される： }d_2^{p,i}.
\end{aligned}
\right.
\tag{4.6.4.2}
\]

\VerdierSemanticAnchor{ch4d-4-6-4-h1}
$H_1^p(X^{\bullet\bullet})$ を次の単純複体と記す：

\VerdierSemanticAnchor{ch4d-4-6-4-3}
\[
\left\{
\begin{aligned}
H_1^p(X^{\bullet\bullet})^i
  &=Z_1^p(X^{\bullet\bullet})^i/B_1^p(X^{\bullet\bullet})^i,\\
d_{H_1^p(X^{\bullet\bullet})}^{i}
  &\text{次の微分から誘導される： }d_2^{p,i}.
\end{aligned}
\right.
\tag{4.6.4.3}
\]

\VerdierSemanticAnchor{ch4d-4-6-4-simple-as-double}
最後に、$Y$ が $\mathcal A$ の単純複体であるとき、第二次数が零でない成分は
零であり、第二次数が零である成分と第一微分は $Y$ の成分と微分であるような
二重複体も $Y$ と記す。

\VerdierSemanticAnchor{ch4d-4-6-5}
\medskip\noindent\textbf{定義 4.6.5.}\quad
$Y$ を $\mathcal A$ の単純複体とする。$Y$ の
\emph{Cartan--Eilenberg 入射分解}とは、二重複体

\VerdierSemanticAnchor{ch4d-4-6-5-double-complex}
\[
I^{\bullet\bullet}\colon\quad
\cdots\longrightarrow0\longrightarrow I^{\bullet,0}
\longrightarrow I^{\bullet,1}\longrightarrow I^{\bullet,2}
\longrightarrow\cdots,
\]
と、\emph{増大写像}と呼ばれる射

\VerdierSemanticAnchor{ch4d-4-6-5-augmentation}
\[
\varepsilon\colon Y\longrightarrow I^{\bullet\bullet}
\]
であって、各整数 $p$ に対し、単純複体
$B_1^p(I^{\bullet\bullet})$ (4.6.4.1) および
$H_1^p(I^{\bullet\bullet})$ (4.6.4.3) は、それぞれ増大写像

\VerdierSemanticAnchor{ch4d-4-6-5-augmentations}
\[
B_1^p(\varepsilon)\colon
\operatorname{Im}(d_Y^{p-1})\longrightarrow B_1^p(I^{\bullet\bullet}),
\]
\[
H_1^p(\varepsilon)\colon
H^p(Y)\longrightarrow H_1^p(I^{\bullet\bullet})
\]
\endgroup
